% Handout for my guest lecture on LaTeX
% B.Tech in Artificial Intelligence
% Global College of Social Science & Technology, Kathmandu
% Compiles with pdflatex, locally or on Overleaf.

\documentclass[10pt,a4paper]{article}

\usepackage[T1]{fontenc}
\usepackage{lmodern}
\usepackage[margin=1.4cm,top=1.0cm,bottom=0.9cm]{geometry}
\usepackage{multicol}
\usepackage{booktabs}
\usepackage{parskip}
\usepackage{xcolor}
\usepackage{hyperref}
\usepackage{qrcode}

\setlength{\parskip}{4pt plus 1pt}

\definecolor{accent}{HTML}{1F6FB2}
\hypersetup{colorlinks=true, urlcolor=accent, linkcolor=accent}

\pagestyle{empty}
\setlength{\columnsep}{0.9cm}

% TODO set the real lecture date before printing
\newcommand{\lecturedate}{\today}

% lrbox lets verbatim text live inside a colorbox
\newsavebox{\codesample}

% Compact colored heading, so the handout fits on one page
\newcommand{\heading}[1]{%
  \vspace{0.3em}
  \noindent{\color{accent}\large\bfseries #1}\par
  \nopagebreak
  \vspace{-0.55em}\noindent{\color{accent}\rule{\linewidth}{0.8pt}}\par
  \vspace{0.1em}
}

\begin{document}

% ---------- Title block with QR code ----------
\noindent
\begin{minipage}[c]{0.70\textwidth}
  {\huge\bfseries Introduction to \LaTeX}\\[5pt]
  {\large Write documents like a scientist}\\[8pt]
  Guest lecture for B.Tech in Artificial Intelligence\\
  Global College of Social Science \& Technology, Kathmandu\\[3pt]
  Danish Puri \quad $\cdot$ \quad \lecturedate
\end{minipage}
\hfill
\begin{minipage}[c]{0.26\textwidth}
  \centering
  \qrcode[nolinks,height=2.5cm]{https://github.com/danish-puri/LaTeX}\\[5pt]
  {\small Slides, notes, exercises}\\
  {\small\url{github.com/danish-puri/LaTeX}}
\end{minipage}

\vspace{0.3em}
\noindent{\color{accent}\rule{\textwidth}{1.2pt}}

\begin{multicols}{2}

\heading{Try it in five minutes}
\begin{enumerate}
  \setlength{\itemsep}{1pt}
  \item Go to \url{overleaf.com} and register. The free account is all you need.
  \item Click \textbf{New Project}, then \textbf{Blank Project}.
  \item Replace the body with the skeleton below.
  \item Press \textbf{Recompile} (Ctrl+Enter, or Cmd+Enter on a Mac).
  \item The PDF appears on the right. Edit, recompile, repeat. That loop is the whole workflow.
\end{enumerate}

\begin{lrbox}{\codesample}
\begin{minipage}{0.95\linewidth}
\begin{verbatim}
\documentclass{article}
\begin{document}
Hello, \LaTeX! Math sits right
in a sentence, $e^{i\pi}+1=0$.
\end{document}
\end{verbatim}
\end{minipage}
\end{lrbox}
\noindent\colorbox{accent!8}{\usebox{\codesample}}

Watch out. The characters \verb|% & _ # $| have special jobs. To print one, escape it, so \verb|\%| \verb|\&| \verb|\_| \verb|\#| \verb|\$|. A bare \verb|%| starts a comment. \LaTeX{} ignores the rest of that line, which is also how you leave notes to yourself.

\heading{Math mode essentials}
Wrap math in \verb|$...$| inside a sentence, or in \verb|\[...\]| to give it its own line.

\begin{center}
\small
\renewcommand{\arraystretch}{1.15}
\begin{tabular}{@{}ll@{}}
\toprule
\textbf{Type this} & \textbf{Get this} \\
\midrule
\verb|$x^2$| & $x^2$ \\
\verb|$a_{ij}$| & $a_{ij}$ \\
\verb|$\frac{a}{b}$| & $\frac{a}{b}$ \\
\verb|$\sqrt{x+y}$| & $\sqrt{x+y}$ \\
\verb|$\alpha, \theta, \pi$| & $\alpha, \theta, \pi$ \\
\verb|$\sum_{i=1}^{n} i$| & $\sum_{i=1}^{n} i$ \\
\verb|$\int_0^1 x \, dx$| & $\int_0^1 x \, dx$ \\
\verb|$\lim_{x \to 0} f(x)$| & $\lim_{x \to 0} f(x)$ \\
\bottomrule
\end{tabular}
\end{center}
\vspace{-0.4em}
{\small The \verb|\,| in the integral row is a thin space.}

You already know enough to typeset real machine learning. Gradient descent is one line.
\[\theta_{t+1} = \theta_t - \eta \, \nabla L(\theta_t)\]
So is softmax.
\[\mathrm{softmax}(z_i) = \frac{e^{z_i}}{\sum_j e^{z_j}}\]

\vfill\columnbreak

\heading{Everyday commands}
\begin{center}
\small
\begin{tabular}{@{}ll@{}}
\toprule
\textbf{Command} & \textbf{What it does} \\
\midrule
\verb|\usepackage{...}| & loads an add-on package \\
\verb|\section{...}| & numbered heading \\
\verb|\textbf{...}| & \textbf{bold} \\
\verb|\textit{...}| & \textit{italics} \\
\verb|itemize| env & bullet list \\
\verb|enumerate| env & numbered list \\
\verb|\includegraphics| & insert an image (needs \verb|graphicx|) \\
\verb|\label| + \verb|\ref| & numbers that fix themselves \\
\verb|\tableofcontents| & builds itself \\
\verb|\cite{...}| & citations via BibTeX \\
\bottomrule
\end{tabular}
\end{center}

\heading{When it breaks}
\begin{itemize}
  \setlength{\itemsep}{1pt}
  \item It will break. That is normal, not a sign you are bad at this.
  \item Read the \textbf{first} error in the log, not the last. Errors cascade. The line number points at or just after the real problem.
  \item Compile often. Then the bug is always in the last thing you typed.
\end{itemize}

The three errors you will meet first, translated.

\begin{center}
\footnotesize
\begin{tabular}{@{}ll@{}}
\toprule
\textbf{The log says} & \textbf{It means} \\
\midrule
\verb|Undefined control sequence| & misspelled command \\
\verb|Missing $ inserted| & math outside \verb|$...$| \\
\verb|\begin{...} ended by \end{...}| & mismatched environment \\
\bottomrule
\end{tabular}
\end{center}

\heading{Your challenge}
Pick one page of any course handout, lab report, or your own notes, and rebuild it in \LaTeX{}. The tables on this page cover everything you need. Stuck on something? Open an issue on the repo and I will answer.

\heading{Keep learning}
\begin{itemize}
  \setlength{\itemsep}{1pt}
  \item Scan the QR code above. The repo has my slides, lecture notes, homework, and a cheat sheet.
  \item \url{overleaf.com/learn} has the best free tutorials.
  \item Need a full paper, thesis, or resume? Never start from scratch. Steal a template from \href{https://www.overleaf.com/gallery}{the Overleaf gallery}.
  \item AI assistants write decent \LaTeX{}. Let them type, but read what they produce and fix it. You are the editor.
\end{itemize}

\end{multicols}

\noindent{\color{accent}\rule{\textwidth}{0.8pt}}\\
{\small I wrote this handout in \LaTeX{}. The source file is in the repo, so once you can read it, you can rebuild this page yourself.}

\end{document}
